\section{The Gaussian time average in Hastings' approximate ground-state projector}
\label{sec:hastings_gaussian_time_average}

% Hastings, arXiv:0705.2024v4, law16.tex, lines 843--864 (step five of the
% proof of Lemma 1). Finite complex matrices; no lattice content.

Step five of the proof of \cite[Lemma~1]{Hastings2007AreaLaw} approximates the
ground-state projector $|\Omega\rangle\langle\Omega|$ by a Gaussian time average
$P_q(H)$ and bounds the perturbation $\|P_q(M)-P_q(H)\|$.

\begin{definition}[Gaussian time average]
    \label{def:hastings_gaussian_time_average}
    \lean{TNLean.MPS.AreaLaw.Hastings.timeAverageVariance,
        TNLean.MPS.AreaLaw.Hastings.gaussianTimeAverage,
        TNLean.MPS.AreaLaw.Hastings.gaussianTimeAverage_eq_integral,
        TNLean.MPS.AreaLaw.Hastings.le_timeAverageVariance_mul_sq,
        TNLean.MPS.AreaLaw.Hastings.sqrt_timeAverageVariance}
    \leanok
    For $q,\Delta E\in\R$, set $h=\max(q/\Delta E^2,0)$ when $\Delta E\ne0$ and
    $h=0$ when $\Delta E=0$. Let $\gamma_h$ be the centered Gaussian measure of
    variance $h$ on $\R$ ($\gamma_0=\delta_0$). For a square matrix $H$, define
    \begin{align}
        P_q(H)=\int e^{iHt}\,d\gamma_h(t)
        \label{eq:hastings_gaussian_time_average}
    \end{align}
    (for $q,\Delta E>0$,
    $P_q(H)=\frac{\Delta E}{\sqrt{2\pi q}}\int e^{iHt-(t\Delta E)^2/(2q)}\,dt$
    as in \cite[lines 843--847]{Hastings2007AreaLaw}). Then $q\le h\,\Delta E^2$
    for $\Delta E\ne0$, and $\sqrt h=\sqrt{\max(q,0)}/\Delta E$ for $\Delta E>0$.
\end{definition}

\begin{theorem}[Gaussian time average across the gap]
    \label{thm:hastings_gaussian_time_average_gap}
    \lean{TNLean.MPS.AreaLaw.Hastings.norm_gaussianTimeAverage_sub_groundProjector_le}
    \leanok
    \uses{def:hastings_gaussian_time_average}
    If $H=H^\dagger$, $\|\Omega\|=1$, $H\Omega=0$, $\Delta E>0$, and
    $H\ge\Delta E\,(1-|\Omega\rangle\langle\Omega|)$, then for every $q\in\R$,
    \begin{align}
        \bigl\|P_q(H)-|\Omega\rangle\langle\Omega|\bigr\|\le e^{-q/2}.
        \label{eq:hastings_gaussian_time_average_gap}
    \end{align}
    This is the estimate of $\|P_q-P_0\|$ at
    \cite[line 848]{Hastings2007AreaLaw}; the source's $q=(l/3)\Delta E/(2v)$
    yields $O(e^{-l/(6\xi)})$ whereas $q=(l/3)\Delta E/v$ recovers the printed
    $O(e^{-l/(3\xi)})$~\cite{gap:hastings07_gaussian_average_exponent}.
\end{theorem}
\begin{proof}\leanok
    \uses{def:hastings_gaussian_time_average}
    In an orthonormal eigenbasis $(e_k)$ of $H$ with $He_k=\lambda_ke_k$
    ($e_0=\Omega$, $\lambda_0=0$, and $\lambda_k\ge\Delta E$ for $k\ne0$),
    (\ref{eq:hastings_gaussian_time_average}) gives
    $\langle e_k,P_q(H)x\rangle=e^{-h\lambda_k^2/2}\langle e_k,x\rangle$, so
    \begin{align}
        \bigl\|P_q(H)x-\langle\Omega,x\rangle\Omega\bigr\|^2
        =\sum_{k\ne0}e^{-h\lambda_k^2}|\langle e_k,x\rangle|^2
        \le e^{-h\Delta E^2}\|x\|^2\le e^{-q}\|x\|^2.
        \notag
    \end{align}
\end{proof}

\begin{theorem}[Comparison of two unitary evolutions]
    \label{thm:hastings_unitary_path_difference}
    \lean{Matrix.norm_hermitianUnitaryPath_sub_le}
    \leanok
    For $M=M^\dagger$, $H=H^\dagger$, and $t\in\R$,
    \begin{align}
        \|e^{iMt}-e^{iHt}\|\le\|M-H\|\,|t|.
        \label{eq:hastings_unitary_path_difference}
    \end{align}
\end{theorem}
\begin{proof}\leanok
    By Duhamel's formula,
    $e^{iHt}-e^{iMt}=e^{iMt}\int_0^te^{-iMs}i(H-M)e^{iHs}\,ds$, and unitarity
    of $e^{iMt},e^{-iMs},e^{iHs}$ gives~(\ref{eq:hastings_unitary_path_difference}).
\end{proof}

\begin{theorem}[First absolute moment of a Gaussian tail]
    \label{thm:hastings_gaussian_absolute_moment_tail}
    \lean{TNLean.Lattice.integral_compl_Icc_abs_gaussianReal_le}
    \leanok
    For $h,T\ge0$ (with the right-hand side $0$ when $h=0$),
    \begin{align}
        \int_{|t|>T}|t|\,d\gamma_h(t)
        \le \sqrt{2h/\pi}\,e^{-T^2/(2h)}.
        \label{eq:hastings_gaussian_absolute_moment_tail}
    \end{align}
\end{theorem}
\begin{proof}\leanok
    For $h=0$, $\gamma_0=\delta_0$ gives $0$. For $h>0$, with
    $\varphi_h(t)=(2\pi h)^{-1/2}e^{-t^2/(2h)}$ and $(-h\varphi_h)'=t\varphi_h$,
    \begin{align}
        \int_{|t|>T}|t|\,d\gamma_h(t)=2\int_T^\infty t\varphi_h(t)\,dt
        =2h\varphi_h(T)=\sqrt{2h/\pi}\,e^{-T^2/(2h)}.
        \notag
    \end{align}
\end{proof}

\begin{theorem}[First absolute moment of a Gaussian]
    \label{thm:hastings_gaussian_absolute_moment}
    \lean{TNLean.Lattice.integral_abs_gaussianReal_le_sqrt_two_mul_div_pi,
        TNLean.Lattice.integral_abs_gaussianReal_le}
    \leanok
    For every $h\ge0$,
    $\int|t|\,d\gamma_h(t)\le\sqrt{2h/\pi}\le\sqrt h$.
\end{theorem}
\begin{proof}\leanok
    \uses{thm:hastings_gaussian_absolute_moment_tail}
    Set $T=0$ in~(\ref{eq:hastings_gaussian_absolute_moment_tail}) and use
    $2\le\pi$.
\end{proof}

\begin{theorem}[Stability of the Gaussian time average]
    \label{thm:hastings_gaussian_time_average_stability}
    \lean{TNLean.MPS.AreaLaw.Hastings.norm_gaussianTimeAverage_sub_le,
        TNLean.MPS.AreaLaw.Hastings.norm_gaussianTimeAverage_sub_groundProjector_le_add}
    \leanok
    \uses{def:hastings_gaussian_time_average,
        thm:hastings_gaussian_time_average_gap}
    For $M=M^\dagger$, $H=H^\dagger$, $\Delta E>0$, and $q\in\R$ (writing
    $\sqrt q=\sqrt{\max(q,0)}$),
    \begin{align}
        \|P_q(M)-P_q(H)\|\le\frac{\sqrt q}{\Delta E}\,\|M-H\|,
        \label{eq:hastings_gaussian_stability_diff}
    \end{align}
    and under the hypotheses of Theorem~\ref{thm:hastings_gaussian_time_average_gap},
    \begin{align}
        \bigl\|P_q(M)-|\Omega\rangle\langle\Omega|\bigr\|
        \le e^{-q/2}+\frac{\sqrt q}{\Delta E}\,\|M-H\|.
        \label{eq:hastings_gaussian_stability_proj}
    \end{align}
    With $M=M_L+M_B+M_R$, (\ref{eq:hastings_gaussian_stability_proj}) is the
    first inequality of \cite[lines 849--864]{Hastings2007AreaLaw}.
\end{theorem}
\begin{proof}\leanok
    \uses{thm:hastings_unitary_path_difference,
        thm:hastings_gaussian_absolute_moment,
        thm:hastings_gaussian_time_average_gap}
    Integrating~(\ref{eq:hastings_unitary_path_difference}) against $\gamma_h$
    and applying Theorem~\ref{thm:hastings_gaussian_absolute_moment} gives
    $\|P_q(M)-P_q(H)\|\le\|M-H\|\int|t|\,d\gamma_h\le(\sqrt q/\Delta E)\|M-H\|$,
    and combining~(\ref{eq:hastings_gaussian_stability_diff}) with
    (\ref{eq:hastings_gaussian_time_average_gap}) gives
    (\ref{eq:hastings_gaussian_stability_proj}).
\end{proof}
